\section{Conditional Gaussian rigidity}\label{sec:rigidity}

The entropy inequality of Section~\ref{sec:entropy} contains the
covariance of a Gaussian tree resolvent.  At height zero that
covariance is singular on a star, because it satisfies the same
eigenvalue equation as the marks.  We first compute its precision
on the constraint space.  The computation has two uses: it fixes
the entropy constants when the height tends to zero, and it splits
each edge precision into positive contributions from its two
endpoints.  Re-rooting will combine these contributions in a
Fisher-information comparison.

This follows the Gaussian-comparison and heat-flow architecture of
Backhausz and Szegedy \cite[Theorem~10.3, Proposition~11.1,
Theorem~11.2, and Lemmas~11.3--11.4]{BackhauszSzegedy2019}.
Here the covariances depend on the entire random environment.
Accordingly, the precision shares are operators depending on that
environment, and their cancellation takes place after expectation
and re-rooting.

\subsection{The Gaussian geometry of a star}

All Green functions in this subsection are boundary values at
\(E+i0\).  At a root \(o\), label the neighbors \(1,\ldots,d\) and
write the incoming cavity messages as
\[
 g_i=c_i+i b_i,\qquad b_i>0,\qquad
 B_*=\sum_{i=1}^d b_i,\qquad
 t_0=V_o-E-\sum_{i=1}^d c_i.
\]
Thus \(m_o=(t_0-iB_*)^{-1}\).  The unsubscripted \(b\) continues
to mean \(d-1\).  Let \(K_o\) be the imaginary part of the tree
Green function compressed to the star, and let \(K_{oi}\) be its
two-coordinate compression to \((o,i)\).  These are spatial
covariances for one color.

\begin{lemma}[Star and edge precision]\label{lem:tree-precision}
The centered Gaussian with covariance \(K_o\) is supported on the
\(d\)-dimensional hyperplane
\begin{equation}\label{eq:star-constraint}
 C:=(V_o-E)Z+\sum_{i=1}^dX_i=0.
\end{equation}
In the coordinate volume
\(\delta(C)\,dZ\prod_i dX_i\), equivalently the volume obtained by
forgetting any one neighbor, its precision quadratic form is
\begin{equation}\label{eq:star-precision}
 \sum_{i=1}^d\frac{|X_i+g_iZ|^2}{b_i}.
\end{equation}
Write \(K_o^{\rm ret}\) for its covariance in coordinates forgetting
neighbor \(d\).  Then
\begin{equation}\label{eq:star-determinant}
 \det K_o^{\rm ret}=|m_o|^2\prod_{i=1}^d b_i,
 \qquad
 h_o^{\rm G}:=\mathcal H(N(0,K_o^{\rm ret}))
 =\frac d2\log(2\pi\mathrm e)+\frac12\sum_i\log b_i+\log|m_o|.
\end{equation}
If \(g'\) is the cavity message from the \(o\)-side of the edge
\((o,i)\), its edge precision and determinant are
\begin{align}
 (Z,X_i)K_{oi}^{-1}(Z,X_i)^{\mathsf T}
 &=\frac{|X_i+g_iZ|^2}{b_i}
       +\frac{|Z+g'X_i|^2}{\Im g'},\label{eq:edge-precision}\\
 \det K_{oi}&=\frac{b_i\Im g'}{|1-g_i g'|^2}.
 \label{eq:edge-determinant}
\end{align}
In particular each edge covariance is positive definite.
\end{lemma}

\begin{proof}
The inverse star Green matrix has central diagonal \(V_o-E\),
root-neighbor entries \(1\), and neighbor diagonals \(1/g_i\).
Its negative imaginary part is diagonal on the neighbors, with
entries \(b_i/|g_i|^2\).  The real central row implies that the
imaginary covariance annihilates the constraint functional in
\eqref{eq:star-constraint}.

To verify its law and its volume normalization directly, start
with independent real Gaussians
\[
 Y_i\sim N(0,b_i),\qquad Z\sim N(0,B_*^{-1}),
\]
and condition on \(t_0Z+\sum_iY_i=0\).  Put \(X_i=Y_i-c_iZ\).
The quadratic form of the resulting density is exactly
\eqref{eq:star-precision}.  Gaussian conditioning gives the
imaginary parts of the Green entries
\[
 m_o,\qquad -m_o g_i,\qquad
 \delta_{ij}g_i+g_i m_o g_j,
\]
in the root, root-neighbor, and neighbor-neighbor positions,
respectively.  Hence this conditioned Gaussian has covariance
\(K_o\).  Integrating its density against \(\delta(C)\) gives
the normalization
\[
 (2\pi)^{d/2}\Bigl(\prod_i b_i\Bigr)^{1/2}|m_o|,
\]
because \(\Var(t_0Z+\sum_iY_i)=(t_0^2+B_*^2)/B_*\).
The coefficient of the forgotten neighbor in \(C\) is one, so
this is precisely the normalization in retained-coordinate
volume, not Euclidean surface volume.  It proves
\eqref{eq:star-determinant}.

On an edge the inverse Green matrix has diagonal entries
\(1/g'\) and \(1/g_i\), and off-diagonal entries \(1\).
Using its diagonal negative imaginary part, or multiplying the
two-by-two matrices, gives \eqref{eq:edge-precision} and
\eqref{eq:edge-determinant}.  Each of the two precision summands
is positive definite: for \(g=c+i\beta\), its matrix is
\begin{equation}\label{eq:one-sided-precision}
 Q(g)=\frac1\beta
 \begin{pmatrix}|g|^2&c\\c&1\end{pmatrix},
 \qquad \det Q(g)=1.
\end{equation}
This also proves edge nondegeneracy.
\end{proof}

We now put the Gaussian reference laws in common Euclidean
coordinates.  Let \(B_i:\RR^d\to\RR^2\) read the edge \((Z,X_i)\)
from the retained star coordinates, reconstructing \(X_d\) from
\eqref{eq:star-constraint} when necessary.  Write
\(U_o^{\rm ret}=(U_o,U_1,\ldots,U_{d-1})\in\RR^{dm}\) for the
retained \(m\)-color star marks, and define
\begin{equation}\label{eq:whitened-star-edge}
 L_o=((K_o^{\rm ret})^{-1/2}\otimes\Id_m)U_o^{\rm ret},
 \qquad
 R_i=\bigl(K_{oi}^{-1/2}B_i(K_o^{\rm ret})^{1/2}\bigr)\otimes\Id_m.
\end{equation}
Then \(R_iR_i^{\mathsf T}=\Id_{2m}\), so \(R_iL_o\) is the
whitened edge mark and \(P_i=R_i^{\mathsf T}R_i\) is the
orthogonal projection onto its reading subspace.

The two precision formulas produce positive matrices
\[
 A_{oi}=\bigl(K_{oi}^{1/2}Q(g_i)K_{oi}^{1/2}\bigr)\otimes\Id_m
\]
on the whitened edge coordinates.  The opposite endpoint defines
\(A_{io}\) using its summand of \eqref{eq:edge-precision}; read
both in the same oriented edge coordinates.  Equations
\eqref{eq:star-precision} and \eqref{eq:edge-precision} give the
pointwise identities
\begin{equation}\label{eq:precision-split}
 \sum_{i=1}^dR_i^{\mathsf T}A_{oi}R_i=\Id_{dm},
 \qquad A_{oi}+A_{io}=\Id_{2m},
 \qquad 0<A_{oi}<\Id_{2m}.
\end{equation}
Thus every star's precision is the sum of its incident
contributions, and the two contributions on an edge give its
whole precision.  These identities are valid in each environment;
no averaging of the covariance has been used.  Other orthogonal
choices of whitened coordinates merely conjugate these formulas.

\subsection{The constant in the entropy comparison}

Define the logarithmic potential of the tree density of states by
\[
 U(\zeta)=\int\log|t-\zeta|\,d\nu_{d,w}(t),\qquad \Im\zeta>0.
\]
The following identity relates it to the local Gaussian volumes.

\begin{lemma}[Tree logarithmic potential]\label{lem:tree-log-potential}
At positive height, with all tree Green functions evaluated at
\(\zeta\),
\begin{equation}\label{eq:tree-log-potential}
 -U(\zeta)=\EE_\omega\log|m_o|
       +\frac d2\EE_\omega\log|1-g_i g'|.
\end{equation}
Both sides have finite limits as \(\zeta=E+i\eta\),
\(\eta\downarrow0\).  Denote the limit of \(U\) by \(U(E)\).
With \(h_{oi}^{\rm G}=\mathcal H(N(0,K_{oi}))\),
\begin{equation}\label{eq:gaussian-entropy-constant}
 \EE_\omega\left[h_o^{\rm G}-\frac12\sum_{i=1}^d h_{oi}^{\rm G}\right]
 =-U(E).
\end{equation}
\end{lemma}

\begin{proof}
Use real parts of analytic logarithms in the upper half-plane.
The edge factor is nonzero, as follows also from its invertible
Schur matrix.  Differentiating the full-root recursion yields
\[
 \partial_\zeta\log m_o
   =m_o\left(1+\sum_i\partial_\zeta g_i\right).
\]
In the expected derivative of the edge term, edge reversal pairs
the two differentiated cavity messages.  Its contribution is
\[
 -d\,\EE_\omega\frac{(\partial_\zeta g_i)g'}{1-g_i g'}
 =-\EE_\omega m_o\sum_i\partial_\zeta g_i,
\]
where \(m_o=g'/(1-g_i g')\).  The indirect derivatives therefore
cancel, leaving \(\EE_\omega m_o\), the analytic derivative whose
real part differentiates \(-U\).  At fixed positive height the
resolvent bounds justify differentiation and expectation.  Both
sides of \eqref{eq:tree-log-potential} are asymptotic to
\(-\log|\zeta|\) at imaginary infinity, which fixes the constant.

The boundary continuity and logarithmic integrability in
Lemma~\ref{lem:tree-log}, together with the Schur identities,
pass the right side to the boundary in expectation.  They apply
also to the two incoming messages in an edge factor, and prove
finiteness of the limit.  Finally insert
\eqref{eq:star-determinant} and \eqref{eq:edge-determinant} into
the Gaussian entropy formula.  The \(\log(2\pi\mathrm e)\)
terms cancel by dimension.  The \(\log b_i\) terms cancel in
expectation because the two cavity messages on an edge have the
same marginal law.  The remaining expression is precisely the
boundary value of \eqref{eq:tree-log-potential}, proving
\eqref{eq:gaussian-entropy-constant}.
\end{proof}

Return to the smoothed field \(F^\eta\) of
Proposition~\ref{prop:conditional-entropy}.  At each fixed
environment, change star coordinates to the retained coordinates
and to the residual \(C\) in each color.  This linear change has
absolute Jacobian one.  The marks have zero residual by
\eqref{eq:mark-normalization-equation}.  The noise residual has
covariance at most \(\eta\Id_m\), since spectral calculus on
the tree gives
\[
 (H_{T_d}-E)\Gamma_\eta(H_{T_d}-E)\le\eta\Id.
\]
The entropy chain rule and the Gaussian maximum-entropy bound
therefore show that
\begin{equation}\label{eq:residual-entropy-bound}
 \mathcal H(F^\eta_{\rm star}\mid\omega)
 \le \mathcal H(F^\eta_{\rm ret}\mid\omega)
           +\frac m2\log(2\pi\mathrm e\,\eta).
\end{equation}
Using this upper bound on the left side of
\eqref{eq:conditional-star-edge-entropy} cancels exactly its
divergent term.  What remains has a finite limit.

We give the entropy-continuity details because the full-star
covariance itself becomes singular.  Take a sequence
\(\eta\downarrow0\) along which the spatial covariances converge
almost surely.  In the retained star and edge coordinates their
limits are positive definite.  At a fixed environment, the
convolved densities consequently converge in \(L^1\), are
uniformly bounded, and have uniformly bounded second moments;
the bin-filling and tail argument from
Section~\ref{sec:entropy} gives entropy convergence.  Conditional
mark second moments are finite almost surely.  For expectation
over environments, positive entropy parts are uniformly
integrable by the Gaussian entropy upper bound, the integrable
conditional mark second moments, and the spatial tree moment
bounds.  Negative parts are controlled by the entropy of the
independent Gaussian noise and the uniformly integrable
log-determinants in Lemma~\ref{lem:tree-log}.  Thus these entropy
limits hold also after conditioning and integration.

Subtract \(m\) times the pure-Gaussian identity
\eqref{eq:gaussian-entropy-constant} and use the whitenings
\eqref{eq:whitened-star-edge}.  Write \(\gamma\) for a standard
Gaussian of the dimension appropriate to each occurrence,
independent of the marks.  With \(\mathcal H_\omega\) denoting
entropy under the conditional law at a fixed environment, the
result is
\begin{equation}\label{eq:whitened-entropy-inequality}
 \mathcal D(1):=
 \EE_\omega\left[
 \mathcal H_\omega(L_o+\gamma)
 -\frac12\sum_{i=1}^d
      \mathcal H_\omega(R_iL_o+\gamma)\right]\ge0.
\end{equation}
Indeed entropy differences from the corresponding pure noise
are unchanged by whitening.  The standard Gaussian constants
then cancel because \(dm=(1/2)d(2m)\).  Branch invariance
replaces the single-edge average by the displayed sum.

\subsection{Heat flow and its equality case}

The precision split now gives a bound in the opposite direction
to \eqref{eq:whitened-entropy-inequality}.  Adding more Gaussian
noise makes its left side increase to zero.  Equality forces
each edge projection to retain all the corresponding score
information, which will determine the conditional density.

\begin{proposition}[Gaussian marks at sampled roots]\label{prop:gaussian-marks}
Within almost every empirical law \(\mathcal L\) of
Proposition~\ref{prop:marked-limits}, the conditional root marks
have law
\begin{equation}\label{eq:conditional-gaussian-marks}
 \Law_{\mathcal L}(U_o\mid\omega)
 =N\left(0,\frac{\Im m_o(E+i0)}{\pi\rho}\Id_m\right).
\end{equation}
In particular the colors are conditionally independent.
For any finite collection of independently sampled roots and
any retained finite position list, their environments are iid
and independent of the limiting point process \(\nu\).
Conditional on these environments and \(\nu\), all their marks
are independent centered Gaussians with the variances in
\eqref{eq:conditional-gaussian-marks}.  Positions of multiplicity
greater than one have one separate color per occurrence.
\end{proposition}

\begin{proof}
Fix \(\mathcal L\), and replace the noise in
\eqref{eq:whitened-entropy-inequality} by \(\sqrt s\,\gamma\),
defining \(\mathcal D(s)\) for \(s>0\).  We first prove
\(\mathcal D'(s)\ge0\), keeping the environment dependence
explicit.

For a smooth positive density \(p\), its score is
\(\nabla\log p\), and its Fisher information is
\(\int|\nabla\log p|^2p\).  At fixed \(\omega\), let \(X\) have
the conditional law of \(L_o+\sqrt s\,\gamma\), let \(\rho_o\)
be its score, and let \(\rho_i\) be the score of \(R_iX\).
Differentiating a marginal density gives
\[
 \EE[ R_i\rho_o(X)\mid R_iX]=\rho_i(R_iX).
\]
Thus conditional Jensen and \eqref{eq:precision-split} give
\begin{align}
 I_o(s)&:=\EE|\rho_o(X)|^2
   =\sum_i\EE\bigl|A_{oi}^{1/2}R_i\rho_o(X)\bigr|^2,\notag\\
 I_o(s)-\sum_i\EE\bigl|A_{oi}^{1/2}\rho_i(R_iX)\bigr|^2
 &=\sum_i\EE\bigl|A_{oi}^{1/2}
       (R_i\rho_o(X)-\rho_i(R_iX))\bigr|^2
 =:\Delta_o(s)\ge0.\label{eq:score-deficit}
\end{align}
All expectations on these lines are conditional on \(\omega\).

The conditional edge law, read in common whitened coordinates,
is the same from its two endpoints.  Invariance under re-rooting
therefore pairs the \(A_{oi}\) term with its \(A_{io}\) term
after \(\EE_\omega\).  Their sum is the identity.  Consequently
\begin{equation}\label{eq:rerooted-fisher}
 \EE_\omega\sum_i\EE
       |A_{oi}^{1/2}\rho_i(R_iX)|^2
 =\frac12\EE_\omega\sum_i I_{oi}(s),
\end{equation}
where \(I_{oi}(s)\) is the ordinary Fisher information of the
edge marginal.  De Bruijn's heat identity now yields
\begin{equation}\label{eq:heat-deficit}
 \mathcal D'(s)=\frac12\EE_\omega\Delta_o(s)\ge0
 \quad\text{for almost every }s>0.
\end{equation}
The comparison of joint and marginal scores belongs to the
Fisher-information approach to entropy subadditivity
\cite{Carlen1991}; for fixed projection families, see
\cite[Proposition~3.3 and Equation~(3.15)]{CarlenCordero2009}.
Combining this comparison with the heat identity \cite{Stam1959}
follows the tree entropy argument of
\cite[Lemma~11.4 and Proposition~11.1]{BackhauszSzegedy2019}.

Here are the analytic justifications for this use of the heat
identity.  Gaussian smoothing of a finite-second-moment law
has finite entropy, and its Fisher information in dimension
\(q\) is at most \(q/s\), by conditioning the Gaussian score.
On a compact positive-time interval, a fixed positive mass of
the input in a bounded ball bounds the absolute logarithm of
the convolved density by \(C(1+|x|^2)\).  Gaussian-kernel
derivative bounds and the input second moment then justify
integration of the heat equation against the log density,
first with spatial cutoffs and then without them.  This gives
\(d\mathcal H(X)/ds=I(X)/2\) in integrated form.  Since
\(0<A_{oi}<\Id\), no inverse-weight moments enter
\eqref{eq:score-deficit}; Fisher information is bounded by
dimension over \(s\), also after environment integration.
The entropy bounds used in the preceding subsection justify
that integration on compact \(s\)-intervals.  These are also
the finite-second-moment hypotheses in
\cite[Appendix~A, Equation~(A.1) and Lemmas~A.2--A.3]{BackhauszSzegedy2019}.

As \(s\to\infty\), \(\mathcal D(s)\to0\).  Indeed, divide
each observation by \(\sqrt s\); the dimensional logarithms
cancel, and at each environment its entropy tends to standard
Gaussian entropy.  The noise entropy is a lower bound, and
the Gaussian entropy upper bound gives domination by a
constant times a logarithmic second moment.  Such a bound is
integrable even if the whitened second moments themselves
are not: with \(M(\omega)=\EE_{\mathcal L}
[\|U_o^{\rm ret}\|^2\mid\omega]\),
\[
 \log\bigl(1+\EE[\|L_o\|^2\mid\omega]\bigr)
 \le\log(1+\|(K_o^{\rm ret})^{-1}\|)+\log(1+M(\omega)).
\]
The first term is integrable by the spatial tree estimates,
and the second by \eqref{eq:mark-normalization-equation} and
invariance.  Edge moments obey the same bound.  Dominated
convergence proves the claim.

Together with \(\mathcal D(1)\ge0\), monotonicity and the zero
limit imply \(\mathcal D(s)=0\) for all \(s\ge1\).  Equation
\eqref{eq:heat-deficit} makes every summand of
\eqref{eq:score-deficit} vanish for almost every such \(s\),
almost surely in \(\omega\).  Choose one \(s\) for which this
holds simultaneously at all the countably many re-rootings.
Positive definiteness of every \(A_{oi}\) gives
\[
 R_i\rho_o(X)=\rho_i(R_iX).
\]
The densities are positive and smooth, so this identity holds
everywhere.  Differentiating shows that the symmetric
log-density Hessian commutes with every \(P_i\).

We next determine what these commuting projections permit.  At one
star, they leave the central law free but force the complementary
spatial directions to be Gaussian with a common color covariance.
We will then read one edge from both endpoints to determine the
central law as well.

For one spatial color, let \(w_c\) be the unit direction
reading the central coordinate, standardized with its noise
variance.  The edge-reading space is
\(\operatorname{span}\{w_c,a_i\}\), where \(a_i\) is the unit
neighbor reading after its central component is removed.
In the Gaussian reference field, conditional on its central value,
the residual neighbor covariance is
\begin{equation}\label{eq:residual-neighbor-covariance}
 \diag(b_i)-\frac{(b_i)_i(b_i)_i^{\mathsf T}}{B_*}.
\end{equation}
It follows that the \(a_i\) span the \((d-1)\)-dimensional
orthogonal complement of \(w_c\), and, for \(i\ne j\),
\begin{equation}\label{eq:residual-directions}
 \langle a_i,a_j\rangle
 =-\sqrt{\frac{b_i b_j}{(B_*-b_i)(B_*-b_j)}}\in(-1,0).
\end{equation}
Here strict inequality in magnitude uses \(d\ge3\).
In particular these lines are pairwise distinct and
nonorthogonal.  For \(m\) colors all these spaces are
tensored with \(\RR^m\).

The intersection of the edge-reading spaces is exactly the
central color space \(\operatorname{span}(w_c)\otimes\RR^m\).
A Hessian commuting with all their projections preserves
this intersection and its orthogonal complement.  On the
complement it preserves each
\(\operatorname{span}(a_i)\otimes\RR^m\), so its restriction
there has the form \(a_i\otimes v\mapsto a_i\otimes D_i v\).
Self-adjointness and \(\langle a_i,a_j\rangle\ne0\) imply
\(D_i=D_j\).  Since the \(a_i\) span, the complementary
Hessian is spatial identity tensored with one color matrix.

Vanishing mixed derivatives first split the log density into
a central function and a complementary function.  In any
orthonormal spatial basis of the complement, the latter has
zero mixed blocks and identical diagonal color Hessians.
It therefore splits into a sum of \(d-1\) functions, each of
its own color variable.  Since \(d-1\ge2\), equality of their
Hessians for all independent choices of these variables
forces the common matrix to be constant.  Integrability of
the density makes it negative definite.  The complementary
law is consequently Gaussian, with covariance equal to
spatial identity tensored with a positive color matrix, and
is independent of the central law.

It remains to constrain that central law.  Read a common edge
from its two incident stars.  In the whitened two-dimensional
spatial edge, their standardized endpoint directions are
distinct by \eqref{eq:edge-determinant}.  They are also
nonorthogonal almost surely.  Indeed the unstandardized endpoint
covariance is
\[
 \Im(-m_o g_i)
 =-\frac{t_0b_i+B_*c_i}{t_0^2+B_*^2}.
\]
Condition on the incoming branches and vary the independent
continuously distributed \(V_o\).  Its numerator is affine
in \(V_o\), with nonzero slope \(b_i\), so it vanishes with
probability zero.  A countable union handles all edges.

In coordinates given by the first endpoint direction and
its perpendicular, the edge log-density Hessian has block
form \(\diag(D_c(z),D)\), with \(D\) the constant Gaussian
color block already obtained.  The other endpoint direction
has coordinates \((\cos\theta,\sin\theta)\), with both
entries nonzero.  Commutation with its projection forces
\[
 \cos\theta\sin\theta\,(D_c(z)-D)=0.
\]
Thus the central block too is the same constant \(D\).
The entire smoothed star is Gaussian with covariance spatial
identity tensored with a color covariance, and the color
covariance agrees across every edge.

Dividing characteristic functions by that of the smoothing
Gaussian now proves the same assertion for the unsmoothed
marks, with a possibly singular positive semidefinite color
covariance \(C_{\mathcal L}(\omega)\).  Their means vanish by
Lemma~\ref{lem:conditional-mean}.  Edge agreement makes
\(C_{\mathcal L}(\omega)\) invariant under re-rooting, and
the covariance is a measurable function of \(\omega\) by
disintegration of the fixed law \(\mathcal L\).

The iid environment is ergodic under tree automorphisms.  One
direct proof approximates an invariant event by an event
depending on a finite ball, moves that ball disjointly far
away by an automorphism, and uses independence; approximation
then makes the invariant event independent of itself.
Root-preserving invariance and re-rooting invariance therefore
make \(C_{\mathcal L}\) a deterministic matrix within this
fixed empirical law.  In original coordinates the root
covariance is \((\Im m_o)C_{\mathcal L}\).  Its expectation
is \(\Id_m\) by Proposition~\ref{prop:marked-limits}, so
\[
 (\EE_\omega\Im m_o)C_{\mathcal L}=\Id_m,
 \qquad C_{\mathcal L}=\frac1{\pi\rho}\Id_m.
\]
This proves \eqref{eq:conditional-gaussian-marks} and fixes the
same answer for every admissible empirical law.

Finally, finite vertex samples were conditionally independent
samples from their empirical law before taking limits.  Their
limiting joint law is therefore the corresponding product
conditional on \(\mathcal L\) and the retained positions.
The one-root marked-environment marginal has just been shown
to be deterministic, with iid environment and the Gaussian
conditional kernel \eqref{eq:conditional-gaussian-marks}.
Its product has no remaining dependence on \(\mathcal L\)
or on \(\nu\).  This proves the sampled-root assertion,
including the conditional formulation and multiplicities.
\end{proof}

Only the root assertion of Proposition~\ref{prop:gaussian-marks}
will be used below.  The finite-dimensional arguments above
require only the fixed spatial covariance and logarithmic
estimates from Section~\ref{sec:tree}.  Tensoring with
\(\Id_m\) changes dimensions and multiplies log-determinants
by \(m\); it does not require higher tree moments as the
finite list length varies.  Thus the disorder threshold
chosen there applies to all the lists retained in the
subsequence construction.
